\section{Tata Letak Divisi IT dalam Perusahaan}

PT Inti Utama Solusindo merupakan perusahaan yang seluruh operasionalnya
berbasis teknologi informasi, di mana hampir seluruh proses bisnis didukung
oleh pengembangan dan implementasi sistem digital. Oleh karena itu, divisi IT
memiliki peran sentral dalam menunjang aktivitas perusahaan, baik dari sisi
pengembangan aplikasi maupun pemeliharaan sistem.

\begin{figure}[H]
    \centering
    \includegraphics[width=10cm]{tataletakdivisi.png}
    \caption{Divisi \textit{Front-End} di PT Inti Utama Solusindo}
\end{figure}

Dalam konteks pelaksanaan magang, penempatan yang diberikan adalah pada bagian
\textit{Front-End} (FE) yang berfokus pada pengembangan UI aplikasi. Saat ini,
Tim \textit{Front-End} terdiri dari tiga anggota, yaitu Dimas Aditya Nugroho
selaku \textit{Supervisor}, Muhammad Syaiful Rahman sebagai \textit{Staff}, dan
Aldian Yohanes sebagai \textit{Intern}.

Pada awal periode magang, ruang lingkup pekerjaan Tim \textit{Front-End}
berfokus pada implementasi tampilan Flutter berbasis pendekatan
\textit{Server-Driven UI} (SDUI) atau \textit{Back-End Driven UI}. Dalam
pendekatan ini, struktur dan konfigurasi UI sebagian besar ditentukan oleh sisi
\textit{back-end} dalam bentuk JSON, sehingga Tim \textit{Front-End}
bertanggung jawab melakukan \textit{rendering} tampilan berdasarkan data dan
konfigurasi yang diterima dari API. Hal ini menyebabkan ruang lingkup pekerjaan
\textit{front-end} pada produk Flutter lebih terintegrasi dengan logika sistem
secara keseluruhan dibandingkan dengan pendekatan \textit{front-end}
konvensional.

Mulai pertengahan periode magang, terjadi transisi \textit{tech stack} dari
Flutter ke React (Next.js + TypeScript) untuk pengembangan ulang antarmuka
HealthyOne dan Canvasser. Pada fase ini, peran yang dijalankan masih berfokus
pada \textit{front-end}, yaitu membangun layar-layar baru menggunakan React
berdasarkan desain yang sama. Selanjutnya, mulai pertengahan April 2026, ruang
lingkup pekerjaan berkembang menjadi \textit{Full-Stack}: layar-layar React
yang sudah dibuat dihubungkan dengan \textit{endpoint} pada \textit{back-end}
Toraja (Django/Python), sehingga pekerjaan tidak lagi terbatas pada sisi
tampilan, tetapi juga mencakup pembuatan dan penyesuaian API di sisi
\textit{back-end}.

Perubahan peran dari \textit{front-end} Flutter menjadi \textit{Full-Stack
    React-Toraja} ini didorong oleh transisi \textit{tech stack} perusahaan secara
keseluruhan serta kebutuhan untuk menangani proyek baru yang lebih kompleks,
yaitu Principal Suite. Dalam skema ini, tim diharapkan tidak hanya berfokus
pada sisi tampilan, tetapi juga memiliki pemahaman dan keterlibatan pada sisi
\textit{back-end} guna mendukung pengembangan sistem secara
\textit{end-to-end}.